% П3.4-1 — Высота прямоугольного треугольника, проведённая к гипотенузе
\documentclass[tikz,border=3pt]{standalone}
% Шапка всех рисунков: те же шрифты и стиль, что в книге.
\def\pfFontPath{../../шрифты/}
\usepackage{fontspec}
\usepackage{unicode-math}
\setmainfont{Manrope}[Path=\pfFontPath, Extension=.ttf,
  UprightFont=*-Regular, BoldFont=*-Bold,
  ItalicFont=*-Regular, ItalicFeatures={FakeSlant=0.18},
  BoldItalicFont=*-Bold, BoldItalicFeatures={FakeSlant=0.18}]
\setmathfont{FiraMath-Regular}[Path=\pfFontPath, Extension=.otf]
\usepackage{стиль/рисунки}

\begin{document}
\begin{tikzpicture}[pf, scale=0.9]
  % --- точки (правим только координаты) ---
  \coordinate (A) at (0,0);
  \coordinate (B) at (6,0);
  \coordinate (C) at (2,2.828);   % угол C прямой: C лежит на окружности с диаметром AB
  \coordinate (H) at (2,0);       % основание высоты

  % --- линии по ролям ---
  \draw[accent] (A) -- (H) node[side,midway,below,text=pfPurple] {$b_c$};
  \draw[accent] (H) -- (B) node[side,midway,below,text=pfPurple] {$a_c$};
  \draw[given] (A) -- (C) node[side,midway,above left] {$b$}
                    -- (B) node[side,midway,above right] {$a$}
                    -- cycle;
  \draw[main] (C) -- (H) node[side,midway,right,text=pfBlue] {$h$};

  % --- прямые углы ---
  \pic[rightang] {right angle = A--C--B};
  \pic[rightang] {right angle = C--H--B};

  % --- подписи точек ---
  \pfPoint{A}{A}{below left}
  \pfPoint{B}{B}{below right}
  \pfPoint{C}{C}{above}
  \pfPoint{H}{H}{below=4pt}
\end{tikzpicture}
\end{document}
